\documentclass[11pt]{amsart}

\usepackage[T1]{fontenc}
\usepackage{lmodern}
\usepackage{microtype}
\usepackage{mathtools,amssymb,amsthm}
\usepackage[hidelinks]{hyperref}

\hypersetup{
  pdftitle={A Proof of Tang and Zhang's Order -2 Schoenberg Conjecture, with
the Equality Case}
}

\newtheorem{theorem}{Theorem}
\newtheorem{lemma}[theorem]{Lemma}
\newtheorem{corollary}[theorem]{Corollary}
\theoremstyle{remark}
\newtheorem*{remark}{Remark}
\newtheorem*{conjecture}{Conjecture C1 (Tang--Zhang)}

\newcommand{\CC}{\mathbb{C}}
\newcommand{\RR}{\mathbb{R}}
\newcommand{\Frob}[1]{\lVert #1\rVert_F}
\newcommand{\tr}{\operatorname{tr}}
\newcommand{\diag}{\operatorname{diag}}
\newcommand{\spec}{\operatorname{spec}}

\title[The order $-2$ Schoenberg conjecture, with equality]
{A Proof of Tang and Zhang's Order $-2$ Schoenberg Conjecture,
with the Equality Case}

\date{}

\subjclass[2020]{Primary 30C15; Secondary 15A42, 30C10, 26C10}
\keywords{Schoenberg type inequality, negative order, critical points of a
polynomial, Schur inequality, normal matrix, equality case}

\begin{document}

\begin{abstract}
Let $P(z)=z\prod_{j=1}^{n}(z-z_j)$ with $z_j\in\CC\setminus\{0\}$, and let
$w_1,\ldots,w_n$ be the zeros of $P'$ counted with multiplicity. Conjecture C1
of Tang and Zhang asserts that
$\sum_k|w_k|^{-2}\le 3\sum_j|z_j|^{-2}+\bigl|\sum_j z_j^{-1}\bigr|^{2}$, and
that equality holds whenever the $z_j$ are collinear with the origin; they
prove it for $n=2$. We prove the inequality for every $n$, and show that
equality holds \emph{if and only if} $z_1,\ldots,z_n$ lie on one line through
the origin. Writing $U=\diag(z_1^{-1},\ldots,z_n^{-1})$ and $J$ for the
all-ones matrix, the proof conjugates the authors' matrix model $B=U(I+J)$,
whose spectrum is $\{w_k^{-1}\}$, by the explicit square root
$Q=(I+J)^{1/2}=I+\frac{\sqrt{n+1}-1}{n}J$: the conjectured right-hand side is
exactly $\Frob{QUQ}^{2}$, so Schur's inequality with its equality case gives
both halves at once, equality occurring precisely when $QUQ$ is normal, which
is precisely collinearity.
\end{abstract}

\maketitle

\section{The conjecture}

For a monic polynomial with a simple zero at the origin,
\begin{equation}\label{eq:P}
  P(z)=z\prod_{j=1}^{n}(z-z_j),
  \qquad z_1,\ldots,z_n\in\CC\setminus\{0\}
  \quad\text{(repetitions allowed),}
\end{equation}
the derivative $P'$ has degree exactly $n$ and no zero at the origin; write
$w_1,\ldots,w_n$ for its zeros, \emph{as a multiset}. Quanyu Tang and Teng
Zhang state the following as Conjecture~C1 of \cite{TZ}, in its section
``Schoenberg type inequality of order $-2$''.

\begin{conjecture}
Let $P(z)=z\prod_{j=1}^{n}(z-z_j)$, $z_j\in\CC\setminus\{0\}$, and let
$w_1,\ldots,w_n$ denote the zeros of $P'(z)$. Then
\begin{equation}\label{eq:C1}
  \sum_{k=1}^{n}\frac{1}{|w_k|^{2}}
  \;\le\;
  3\sum_{j=1}^{n}\frac{1}{|z_j|^{2}}
  +\Biggl|\sum_{j=1}^{n}\frac{1}{z_j}\Biggr|^{2}.
\end{equation}
Moreover, equality in \eqref{eq:C1} holds whenever $z_1,\ldots,z_n$ all lie on
the same line passing through the origin.
\end{conjecture}

Tang and Zhang prove \eqref{eq:C1} for $n=2$.

\section{Notation and the result}

Throughout, $n\ge1$ and $z_1,\ldots,z_n$ are as in \eqref{eq:P}. Put
\begin{gather}
  u_j=\frac{1}{z_j},\quad
  U=\diag(u_1,\ldots,u_n),\quad
  S=\sum_{j=1}^{n}u_j,\quad
  D=\diag(z_1,\ldots,z_n),
  \label{eq:notation}\\
  J=\mathbf{1}\mathbf{1}^{T}\ \text{(all ones)},\quad
  M=I+J,\quad
  c=\frac{\sqrt{n+1}-1}{n},\quad
  Q=I+cJ,
  \label{eq:notation2}\\
  K=QUQ,
  \quad\text{that is}\quad
  K_{jk}=\delta_{jk}u_j+c\,(u_j+u_k)+c^{2}S .
  \label{eq:K}
\end{gather}
Say that $z_1,\ldots,z_n$ are \emph{collinear with the origin} if they all lie
on a single line through $0$, i.e.\ if $z_j=\rho r_j$ for some
$\rho\in\CC\setminus\{0\}$ and some real $r_j\ne0$; equivalently, if
$z_j\overline{z_k}\in\RR$ for all $j,k$. For $n=1$ the condition is vacuous.

\begin{theorem}\label{thm:main}
For every $n\ge1$ and all $z_1,\ldots,z_n\in\CC\setminus\{0\}$, repetitions
allowed, the inequality \eqref{eq:C1} holds. Equality holds in \eqref{eq:C1}
if and only if $z_1,\ldots,z_n$ are collinear with the origin.
\end{theorem}

Only the inequality and the ``whenever'' half of the equality clause are
stated in \cite{TZ}; the ``only if'' half does not appear there. We use Schur's
inequality in the form in which it carries its own equality case.

\begin{lemma}[Schur]\label{lem:schur}
For every $X\in\CC^{n\times n}$ with eigenvalues $\lambda_1,\ldots,\lambda_n$
listed with multiplicity, $\sum_k|\lambda_k|^{2}\le\Frob{X}^{2}$, with equality
if and only if $X$ is normal.
\end{lemma}

The inequality is \cite{Schur}; for the equality case see
\cite[\S2.5]{HornJohnson}. It is quoted in this form, with the equality case,
in \cite{TZ}.

\section{Proof of Theorem~\ref{thm:main}}

\subsection*{Step 1: an explicit square root of $M=I+J$}

Since $J^{2}=nJ$,
\[
  Q^{2}=(I+cJ)^{2}=I+(2c+nc^{2})J=I+J=M,
\]
because $nc^{2}+2c-1=0$ by the choice of $c$ in \eqref{eq:notation2}. The
matrix $Q$ is real symmetric with positive entries, and it is invertible with
$Q^{-1}=I+c'J$, $c'=\bigl(\tfrac{1}{\sqrt{n+1}}-1\bigr)/n$, since
$c+c'+ncc'=0$. Note $M^{2}=I+(n+2)J\ne M$: the sandwich matrix of the order
$-2$ problem is not idempotent, which is why a square root is taken here
rather than a projection inserted.

\subsection*{Step 2: the spectrum of $K$, as a multiset}

Let $A=(I-\tfrac{1}{n+1}J)D$. Since $JD$ is the rank-one matrix $\mathbf1z^{T}$,
the Schur determinant formula gives
\[
  \det(wI-A)
  =\prod_{j=1}^{n}(w-z_j)\cdot
   \Bigl(1+\tfrac{1}{n+1}\sum_{j=1}^{n}\tfrac{z_j}{w-z_j}\Bigr).
\]
Substituting $\tfrac{z_j}{w-z_j}=-1+\tfrac{w}{w-z_j}$ turns the bracket into
$\tfrac{1}{n+1}\bigl(1+w\sum_j\tfrac{1}{w-z_j}\bigr)$, and since
$P'(w)=\prod_j(w-z_j)\cdot\bigl(1+w\sum_j\tfrac{1}{w-z_j}\bigr)$ we get the
identity
\begin{equation}\label{eq:charpoly}
  \det(wI-A)=\frac{P'(w)}{n+1},
\end{equation}
whose right-hand side is monic of degree $n$. Hence
$\spec(A)=\{w_1,\ldots,w_n\}$ \emph{as a multiset}, and
$\prod_kw_k=e_n(z)/(n+1)\ne0$, so no $w_k$ vanishes and $A$ is invertible.
Because $(I-\tfrac{1}{n+1}J)(I+J)=I$ exactly, we may read off the inverse:
\begin{equation}\label{eq:B}
  B:=A^{-1}=D^{-1}(I+J)=UM,
  \quad
  \spec(B)=\Bigl\{\tfrac{1}{w_1},\ldots,\tfrac{1}{w_n}\Bigr\}
  \ \text{with mult.}
\end{equation}

Finally $K=QUQ=Q(UM)Q^{-1}=QBQ^{-1}$ is similar to $B$, so
$\spec(K)=\{w_k^{-1}\}$ with multiplicity, and
\begin{equation}\label{eq:lhs}
  \sum_{k=1}^{n}\frac{1}{|w_k|^{2}}
  =\sum_{k=1}^{n}\bigl|\lambda_k(K)\bigr|^{2}.
\end{equation}

The multiset reading in \eqref{eq:charpoly} is what \eqref{eq:C1} requires,
since it sums over the $n$ zeros of $P'$ with multiplicity. At $z=(1,1,1)$ one
has $P'(z)=4z^{3}-9z^{2}+6z-1=(z-1)^{2}(4z-1)$, so $w=1$ is a double zero and
\eqref{eq:C1} closes as $1+1+16=18=9+9$.

\subsection*{Step 3: the conjectured right-hand side is $\Frob{K}^{2}$}

Since $Q$ is real symmetric, $K^{*}=Q\overline UQ$, so using $Q^{2}=M$,
\begin{align*}
  \Frob{K}^{2}
  &=\tr(KK^{*})=\tr(QUQ\cdot Q\overline UQ)
   =\tr(Q^{2}UM\overline U)\\
  &=\tr(MUM\overline U)
   =\sum_{j,m}M_{jm}^{2}\,u_m\overline{u_j}.
\end{align*}
Now $M_{jm}^{2}=(1+\delta_{jm})^{2}=1+3\delta_{jm}$ and
$\sum_{j,m}u_m\overline{u_j}=|S|^{2}$, whence
\begin{equation}\label{eq:rhs}
  \Frob{K}^{2}=3\sum_{j=1}^{n}|u_j|^{2}+|S|^{2}
  =3\sum_{j=1}^{n}\frac{1}{|z_j|^{2}}
   +\Biggl|\sum_{j=1}^{n}\frac{1}{z_j}\Biggr|^{2},
\end{equation}
which is the right-hand side of \eqref{eq:C1} exactly.

\subsection*{Step 4: the inequality}

Combine \eqref{eq:lhs}, \eqref{eq:rhs} and Lemma~\ref{lem:schur} applied to
$X=K$:
\[
  \sum_{k}\frac{1}{|w_k|^{2}}
  =\sum_k|\lambda_k(K)|^{2}
  \le\Frob{K}^{2}
  =3\sum_j\frac{1}{|z_j|^{2}}+\Bigl|\sum_j\frac1{z_j}\Bigr|^{2},
\]
with equality if and only if $K$ is normal. This is \eqref{eq:C1} for every
$n$.

\subsection*{Step 5: $K$ is normal exactly when the $z_j$ are collinear with
the origin}

We have $KK^{*}=Q\,(UM\overline U)\,Q$ and $K^{*}K=Q\,(\overline UMU)\,Q$ with
$Q$ invertible, so $K$ is normal if and only if
$UM\overline U=\overline UMU$. Entrywise, $U$ being diagonal and
$M_{jk}=1+\delta_{jk}$,
\[
  (UM\overline U)_{jk}=(1+\delta_{jk})\,u_j\overline{u_k},
  \qquad
  (\overline UMU)_{jk}=(1+\delta_{jk})\,\overline{u_j}u_k .
\]
The diagonal entries agree always, both being $2|u_j|^{2}$. Off the diagonal
the factor is $1$, so normality says precisely that
$u_j\overline{u_k}\in\RR$ for all $j\ne k$. Writing
$u_j=\rho_je^{i\varphi_j}$ with $\rho_j>0$, this says
$\varphi_j\equiv\varphi_k \pmod \pi$ for all $j,k$, a transitive condition;
that is, all $u_j$ lie on one line through $0$. Since
$1/(re^{i\theta})=r^{-1}e^{-i\theta}$, the $u_j$ lie on a line through $0$ if
and only if the $z_j$ do. Combining with Step~4 proves
Theorem~\ref{thm:main}. \qed

\begin{remark}
The ``if'' half is also visible directly: if $z_j=\rho r_j$ with $r_j$ real,
then $U=\rho^{-1}R$ with $R$ real diagonal, so $K=\rho^{-1}(QRQ)$ is a scalar
times a real symmetric matrix, hence normal.
\end{remark}

\begin{remark}[method]
Steps~3 and~5 transpose a template published earlier by one of the two authors
of \cite{TZ} and cited there, namely \cite{Tang25}, where the projection
$S=I-\tfrac1nJ$ plays the role of $Q$ and the same passage from normality to
collinearity is carried out; the step added here is the use of the square root
$Q=M^{1/2}$, $M=I+J$ not being idempotent.
\end{remark}

\section{Two instances}

Write $L(z)$ and $R(z)$ for the two sides of \eqref{eq:C1}.

\medskip
\noindent\textbf{Equality at $n=3$, $z=(1,2,3)$.}
Here $n+1=4$ is a square, so $c=\tfrac13$ exactly and everything is rational.
$P(z)=z^{4}-6z^{3}+11z^{2}-6z$ and $P'(z)=4z^{3}-18z^{2}+22z-6$, so the
elementary symmetric functions of $w_1,w_2,w_3$ are
$(e_1,e_2,e_3)=(\tfrac92,\tfrac{11}2,\tfrac32)$ and, all $w_k$ being real,
\begin{align*}
  L(z)&=\sum_k\frac{1}{w_k^{2}}
      =\frac{e_2^{2}-2e_1e_3}{e_3^{2}}
      =\frac{\tfrac{121}{4}-\tfrac{27}{2}}{\tfrac94}
      =\frac{67}{9},\\
  R(z)&=3\cdot\frac{49}{36}+\Bigl(\frac{11}{6}\Bigr)^{2}=\frac{67}{9}.
\end{align*}
Through the matrix of \eqref{eq:K}: with $u=(1,\tfrac12,\tfrac13)$,
$S=\tfrac{11}6$ and $c=\tfrac13$,
\[
  K=\frac{1}{54}
  \begin{pmatrix}
    101 & 38 & 35\\
     38 & 56 & 26\\
     35 & 26 & 41
  \end{pmatrix},
  \qquad
  \Frob{K}^{2}=\frac{21708}{2916}=\frac{67}{9},
\]
the squared numerators summing to $21708$. This $K$ is real symmetric, hence
normal, as Step~5 requires.

\medskip
\noindent\textbf{Strict at $n=2$, $z=(1,i)$.}
$u_1\overline{u_2}=i\notin\RR$, so $K$ is not normal. $P'(z)=3z^{2}-2(1+i)z+i$,
and $L=6<8=R$.

\section{Verification}

The algebraic identities used above, and every number displayed in \S4, are
re-derived by the program \texttt{verify.py} accompanying this note, in exact
arithmetic over $\mathbb{Q}(\sqrt{n+1}\,)(i)$ with no floating point anywhere
and no third-party package; its recorded output is \texttt{verify.output.txt}
($67$ checks, all passing). Checked there are the identity $nc^{2}+2c-1=0$ and
$Q^{2}=I+J$ symbolically for $n\le12$; the determinant identity
\eqref{eq:charpoly}, the inverse \eqref{eq:B}, the similarity $KQ=QB$ and the
Frobenius identity \eqref{eq:rhs} on batteries of Gaussian-rational tuples; the
equivalence ``$K$ normal $\iff$ collinear with the origin'' on a battery with
both classes non-empty; and an exact decision of \eqref{eq:C1} at $n=2$ on
$250$ Gaussian-rational pairs, where the inequality holds always and equality
occurs exactly on the collinear pairs. The program does not prove
Lemma~\ref{lem:schur}, and it decides \eqref{eq:C1} only at $n=2$; for
$n\ge3$ both halves of Theorem~\ref{thm:main} rest on the proof of \S3 and not
on any computation. Nothing is claimed here about the other Schoenberg type
inequalities of \cite{TZ}, of orders $-1$, $-4$, $-2m$ and $-p$.

\begin{thebibliography}{9}

\bibitem{HornJohnson}
R.~A. Horn and C.~R. Johnson,
\emph{Matrix Analysis}, 2nd ed.,
Cambridge University Press, Cambridge, 2013.
(Schur's inequality with its equality case: \S2.5.)

\bibitem{Schur}
I.~Schur,
\emph{\"Uber die charakteristischen Wurzeln einer linearen Substitution mit
einer Anwendung auf die Theorie der Integralgleichungen},
Math. Ann. \textbf{66} (1909), no.~4, 488--510.
\href{https://doi.org/10.1007/BF01450045}{doi:10.1007/BF01450045}.

\bibitem{Tang25}
Q.~Tang,
\emph{Schoenberg type inequalities},
arXiv:2504.09837v2, 2025.
\href{https://arxiv.org/abs/2504.09837}{arXiv:2504.09837}.

\bibitem{TZ}
Q.~Tang and T.~Zhang,
\emph{Sharp Schoenberg type inequalities and the de Bruin--Sharma problem},
arXiv:2508.10341v3, 2025.
\href{https://arxiv.org/abs/2508.10341}{arXiv:2508.10341}.
Conjecture~C1 is as numbered in the v3 preprint.

\end{thebibliography}

\end{document}
